\documentclass[12pt]{article}

\usepackage{amsmath, amssymb, amsthm}
\usepackage{enumitem}
\usepackage{hyperref}
\usepackage{geometry}
\geometry{a4paper, margin=1in}

% Define theorem styles
\newtheorem{theorem}{Theorem}
\newtheorem{lemma}[theorem]{Lemma}
\newtheorem{proposition}[theorem]{Proposition}
\newtheorem{corollary}[theorem]{Corollary}
\newtheorem{definition}[theorem]{Definition}

% Document title and author
\title{JAIO II: Homework 3}
\author{Heorhii Lopattn}
\date{}

\begin{document}

\maketitle

%%%%%%%%%%%%%%%%%%%%%%%%%%%%%%%%%%%%%%%%%%%%%%%%%%%%%%%%%%%%%%%%%%%%
\section*{Problem 1}
\textbf{Statement:} Show that the following problem is decidable:  
\begin{itemize}
    \item \textbf{Input:} A deterministic register automaton, defining a language \(L \subseteq \mathcal{A}^*\).
    \item \textbf{Question:} Does the language satisfy, for every word \(w \in \mathcal{A}^*\) and function \(\sigma : \mathcal{A} \to \mathcal{A}\), not necessarily a permutation, 
    \[
    w \in L \iff \sigma(w) \in L?
    \]
\end{itemize}

\subsection*{Solution:}
\begin{proof}
We first show that if 
\[
\forall w\in\mathcal{A}^*,\,\forall\sigma:\mathcal{A}\to\mathcal{A}:\quad w\in L\iff \sigma(w)\in L,
\]
then membership in \(L\) depends only on the length of the word. Indeed, consider any fixed letter \(a\in \mathcal{A}\) and define the constant function \(\sigma(a')=a\) for every \(a'\in\mathcal{A}\). Then for every \(w\in \mathcal{A}^*\) we have \(\sigma(w)=a^{|w|}\). Hence,
\[
w\in L \iff a^{|w|}\in L.
\]
Thus, for every \(n\ge 0\) either all words of length \(n\) belong to \(L\) or none do. 

Now we need to check if all paths of a given length are either not accepting or all accepting. In order to do that: 

\begin{itemize}
\item rename all transitions to be on the same letter, essentially restricting the input alphabet to a single letter. Now we have an NRA instead of DRA.
\item take all possible combinations of states where not all states are accepting and not all states are rejecting. Create a disjoint union construction with as many copies of the NRA as there are combinations, identifying an accepting state in each copy as such that there is a run to all states in the respective combination (using product construction).
\item if in this construction exists an accepting run then the original condition doesn't hold, and if there is no such run, then it holds.
\item from exercises, the emptiness problem is PSPACE, so the original problem is decidable.
\end{itemize}
\end{proof}

\newpage

%%%%%%%%%%%%%%%%%%%%%%%%%%%%%%%%%%%%%%%%%%%%%%%%%%%%%%%%%%%%%%%%%%%%
\section*{Problem 2}
\textbf{Statement:} A language of infinite words \(L \subseteq \Sigma^\omega\) is \textit{closed} if the following condition is satisfied: For every infinite word \(w = a_0a_1\ldots \in \Sigma^\omega\), if every finite prefix \(a_0a_1\ldots a_n \in \Sigma^*\) of \(w\) can be extended to an \(\omega\)-word \(a_0a_1\ldots a_n \cdot v \in L\), then \(w \in L\).

\begin{enumerate}[label=\textbf{\arabic*.}]
    \item Show that there is an \(\omega\)-regular language which is not closed.
    \item Show that there exists a closed language which is not \(\omega\)-regular.
    \item Show that the following problem is decidable: Given a nondeterministic Büchi automaton, decide whether the language it recognizes is closed.
\end{enumerate}

\subsection*{Solution:}
\begin{enumerate}[label=\textbf{\arabic*.}]
    \item 
    \begin{proof}
    Consider the language
    \[
    L = \{ w\in\{a,b\}^\omega \mid w \text{ contains infinitely many occurrences of } a\}.
    \]
    It is known that $L$ is $\omega$-regular, however it is not closed, since in a word $b^\omega$ every finite prefix can be extended to a word in $L$, however the word itself is not in $L$.
    \end{proof}
    
    \item 
    \begin{proof}
      There are countably many $\omega$-regular languages, but uncountably many closed languages, since for each infinite word $w\in\Sigma^\omega$ there is a unique closed language $L_w$ that contains $w$ and no other word. Hence, there exists a closed language that is not $\omega$-regular by counting argument.
          
          
    \end{proof}
    
    \item 
    \begin{proof}
      It would suffice to show that the closure of a $\omega$-language is $\omega$-regular. From that, we would simply check emptiness of the closure of the language, minus the original automata, which is decidable.
  The automaton for closure can be obtained by extending the set of accepting states.
    \end{proof}
\end{enumerate}

\end{document}
